% ===== PRÉAMBULE CANONIQUE — à coller VERBATIM en tête de CHAQUE .tex =====
\documentclass[11pt,a4paper]{article}
\usepackage[utf8]{inputenc}\usepackage[T1]{fontenc}\usepackage{lmodern}
\usepackage[french]{babel}
\usepackage{amsmath,amssymb,mathtools}\usepackage{icomma}
\usepackage{siunitx}\sisetup{locale=FR}  % \qty{}{}, \si{}, \ang{}
\usepackage{xcolor}\usepackage{colortbl}
\usepackage{tikz}\usepackage{pgfplots}\pgfplotsset{compat=1.18}
\usepackage[siunitx,europeanresistors]{circuitikz} % circuits (élec.)
\usepackage[most]{tcolorbox}\usepackage{enumitem}\usepackage{array,booktabs}
\usepackage[a4paper,margin=2.2cm]{geometry}
\usetikzlibrary{arrows.meta,calc,angles,quotes,positioning,patterns,%
  backgrounds,shapes.geometric,decorations.pathmorphing,decorations.markings,intersections}
\definecolor{primary}{RGB}{222,49,99}\definecolor{primarydark}{RGB}{150,22,55}
\definecolor{defcol}{RGB}{0,114,178}\definecolor{methcol}{RGB}{52,92,148}
\definecolor{excol}{RGB}{0,135,90}\definecolor{attncol}{RGB}{200,40,55}
\tcbset{boxbase/.style={enhanced,breakable,boxrule=0.9pt,arc=2.5pt,
  left=3mm,right=3mm,top=2.3mm,bottom=2.3mm,fonttitle=\bfseries,coltitle=white}}
% --- boîtes TUTORIEL (noms EXACTS, ne pas renommer) ---
\newtcolorbox{cle}[1][]{boxbase,colback=defcol!6,colframe=defcol,
  title={\textsc{À retenir}\ifx\relax#1\relax\relax\else\textmd{ : \itshape #1}\fi}}
\newtcolorbox{methode}[1][]{boxbase,colback=methcol!7,colframe=methcol,sharp corners,
  title={\textsc{Méthode}\ifx\relax#1\relax\relax\else\textmd{ --- \itshape #1}\fi}}
\newtcolorbox{exemplebox}[1][]{boxbase,colback=excol!5,colframe=excol,
  title={\textsc{Exemple}\ifx\relax#1\relax\relax\else\textmd{ : \itshape #1}\fi}}
\newtcolorbox{piege}[1][]{boxbase,colback=attncol!6,colframe=attncol,
  title={\textsc{Piège}\ifx\relax#1\relax\relax\else\textmd{ : \itshape #1}\fi}}
\newtcolorbox{remarque}[1][]{boxbase,colback=black!4,colframe=black!45,
  title={\textsc{Remarque}\ifx\relax#1\relax\relax\else\textmd{ : \itshape #1}\fi}}
% --- boîtes EXERCICE / CORRIGÉ (noms EXACTS, auto-comptées) ---
\newtcolorbox[auto counter]{exo}[1][]{boxbase,colback=primary!4,colframe=primary,
  title={\textsc{Exercice}~\thetcbcounter\ifx\relax#1\relax\relax\else\textmd{ --- \itshape #1}\fi}}
\newtcolorbox[auto counter]{corrige}[1][]{boxbase,colback=excol!4,colframe=excol,
  title={\textsc{Corrigé}~\thetcbcounter\ifx\relax#1\relax\relax\else\textmd{ --- \itshape #1}\fi}}
\newcommand{\vv}[1]{\vec{#1}}\newcommand{\ds}{\displaystyle}
\newcommand{\R}{\mathbb{R}}\newcommand{\N}{\mathbb{N}}\newcommand{\Z}{\mathbb{Z}}
% =========================================================================
\begin{document}
\usetikzlibrary{babel}
\begin{exo}[retrouver la longueur]
Un fil de cuivre ($\rho=\qty{1.7e-8}{\ohm\metre}$) de section $S=\qty{1}{\milli\metre\squared}$
présente une résistance $R=\qty{2}{\ohm}$. Quelle est sa longueur $L$ ?
\end{exo}

\begin{exo}[comparer deux fils]
Deux fils sont faits du même matériau. Le fil 2 a une longueur \textbf{3 fois} plus grande et une
section \textbf{2 fois} plus grande que le fil 1. Sachant que $R_1=\qty{4}{\ohm}$, calcule $R_2$ en
raisonnant sur le rapport $\dfrac{R_2}{R_1}$.
\end{exo}

\begin{exo}[la résistance double avec la température]
Un fil d'aluminium a une résistance $R_0=\qty{50}{\ohm}$ à $\qty{0}{\degreeCelsius}$ et un coefficient
$\alpha=\qty{4e-3}{\per\degreeCelsius}$.
\begin{enumerate}[label=\alph*)]
  \item Calcule le rapport $\dfrac{R(T)}{R_0}$ à $T=\qty{250}{\degreeCelsius}$.
  \item En déduire $R$ à $\qty{250}{\degreeCelsius}$.
\end{enumerate}
\end{exo}

\begin{exo}[classer des matériaux]
On dispose de trois matériaux de résistivités : cuivre $\qty{1.7e-8}{\ohm\metre}$, fer
$\qty{10e-8}{\ohm\metre}$, carbone $\qty{3.5e-5}{\ohm\metre}$.
\begin{enumerate}[label=\alph*)]
  \item Classe-les du \emph{meilleur} au \emph{moins bon} conducteur.
  \item Lequel a la plus grande conductivité $\sigma$ ?
\end{enumerate}
\end{exo}

\begin{exo}[identifier un matériau]
Un fil de longueur $L=\qty{100}{\metre}$ et de section $S=\qty{2}{\milli\metre\squared}$ a une
résistance mesurée $R=\qty{0.85}{\ohm}$.
\begin{enumerate}[label=\alph*)]
  \item Calcule la résistivité $\rho$ du matériau.
  \item De quel matériau s'agit-il, d'après la table ? (Ag $1{,}6$ ; Cu $1{,}7$ ; Al $2{,}8$, en
        $10^{-8}\,\si{\ohm\metre}$.)
\end{enumerate}
\end{exo}

\end{document}
